% !TEX root = ../main.tex
\chapter{Introduction}

International remittances are one of the main private channels through which local economic shocks can cross national borders. Mexico is a central setting in which to study this mechanism. Banco de M{\'e}xico records US\$64.7 billion in family remittance receipts in 2024, while World Bank estimates place Mexico among the largest remittance recipients in the world, second only to India \parencite{banxico_remittances_ca11, worldbank_remittances}. The U.S.--Mexico corridor is especially important because Mexican households receive transfers from migrants whose earnings are tied to local labor markets across the United States. This paper asks whether a negative local shock in the United States is transmitted to Mexican municipalities through remittances. Do migrants smooth transfers after being hit by a shock, or do origin households bear part of the adjustment through lower remittance inflows?

The answer matters because remittances are not only an aggregate balance-of-payments flow, they are also part of household resources. Our data do not directly observe Mexican household income or consumption. Even so, changes in remittance inflows are economically meaningful for household welfare. Prior work shows that remittances are an important source of household resources in Mexico, affecting consumption, poverty, and inequality \parencite{adamsetal2005}, food security \parencite{MORARIVERA2021105349}, and schooling and child labor decisions \parencite{ALCARAZ2012156}. We therefore interpret changes in municipal remittance inflows as reduced-form evidence on a channel with direct relevance for household income and consumption, while not claiming to estimate those household outcomes themselves.

The main empirical difficulty is that most large shocks are not cleanly assigned to one side of the migration corridor. \textcite{caballero_2023} show that Mexican households connected to U.S. labor markets more severely affected by the Great Recession were less likely to receive remittances, providing important evidence that migrant networks transmit U.S. shocks to Mexico. Yet the Great Recession simultaneously affected the United States and Mexico, making it difficult to isolate the remittance channel from other aggregate forces. Their analysis also focuses on the extensive margin of remittance receipt. Quantifying how much bilateral remittance flows adjust on the intensive margin is necessary to measure the income loss experienced by origin households.

We address these challenges by exploiting Hurricane Ian, which made landfall in southwestern Florida on September 28, 2022. Ian caused an estimated US\$109.5 billion in damages in Florida and was the costliest hurricane in the state's history \parencite{nhc_ian_2026}. The shock was severe and spatially concentrated on the sender side of the corridor: it disrupted economic activity in Florida but did not directly affect Mexican municipalities. Natural disasters in the United States can reduce local productivity, labor demand, housing wealth, and household resources \parencite{natural_disasters_US}. Hurricane Ian therefore provides a localized negative shock to the capacity of Florida-based migrants to remit, while leaving recipient municipalities in Mexico exposed only through their pre-existing migration links to Florida.

We base our argument on the following idea. A shock to Florida migrants can directly reduce remittances from Florida to Mexican municipalities. However, many origin municipalities receive remittances from migrants in several U.S. states. If migrants care about total resources available to the origin household, then a fall in Florida remittances can raise the marginal value of transfers from migrants living elsewhere. The same shock can therefore generate two effects at once: a direct decline in Florida-to-Mexico remittances and a compensatory response from other U.S. states connected to the same municipalities.

The first part of the paper formalizes this idea with a structural remittance model. The model follows the insight in \textcite{albert_monras_2022} that migrants allocate resources across where they live and where their origin households are located. In our setting, this idea is applied to remittance choices rather than to location choices. Migrants choose between local consumption in the United States and resources available to the origin household in Mexico. They differ in income and in the strength of their altruistic or obligation motive, and remitting involves both a fixed cost and a variable cost. In the spirit of models with heterogeneous agents, fixed costs, and variable costs \parencite{melitz2003}, the fixed cost generates an extensive margin of remitting, while the variable cost governs the intensity of transfers.

The key extension is that origin-household resources include local non-remittance income plus remittances received from migrants in other U.S. states. This makes the model directly useful for interpreting the empirical patterns. A negative shock to Florida migrants reduces their optimal remittances. At the same time, for households linked to Florida, the fall in Florida transfers lowers perceived origin resources for migrants in other states, raising the marginal utility of additional remittances under an insurance motive. The model therefore provides a unified way to interpret both the direct exposure effect and the compensatory network response observed in the data.

The second part of the paper brings the model to the data by constructing a bilateral panel of remittance flows from U.S. states to Mexican municipalities. Because such flows are not directly observed, we combine two sources from Banco de M{\'e}xico with MCAS administrative records. Banco de M{\'e}xico reports remittance inflows by U.S. state of origin and by Mexican municipality of receipt \parencite{banxico_ce168, banxico_ce166}. MCAS records identify the Mexican municipality of origin and U.S. state of residence of Mexican consular-card holders, providing a proxy for the geography of migrant networks. We use these records to construct state-to-municipality migration shares and allocate observed state-level remittance outflows across Mexican municipalities.

Two validation exercises support this construction. First, the migration network data are highly persistent: the distribution of Mexican migrants across U.S. states and Mexican municipalities has year-to-year correlations above 0.90 from 2012 onward. Second, when we use these migration shares to allocate state-level remittance outflows to Mexican municipalities, the resulting municipal aggregates correlate between 0.75 and 0.85 with official municipal remittance receipts. These facts suggest that the estimated corridors capture meaningful variation in the geography of remittances.

The third part of the paper studies the shock in three steps. First, we examine remittance outflows from Florida after Hurricane Ian and observe a decrease in the data, consistent with a sender-side shock reducing migrants' capacity to remit. Second, we exploit variation in Mexican municipalities' pre-existing exposure to Florida to estimate whether Ian affected municipal remittance inflows differently depending on their exposure to the shocked corridor. In the period immediately following Ian, a 1\% increase in a municipality's pre-existing exposure to Florida is associated with a 0.16\% decrease in remittances received. This result is consistent with a sender-side income shock being transmitted through the remittance corridor to origin municipalities.

Finally, we examine migration-network spillovers. Among municipalities jointly exposed to Florida and to other large sender states, remittances from California and Texas rise relative to Florida in the first post-Ian quarters. The diagnostic plots show short-run increases of roughly 20--40 percent for these alternative corridors before the response fades and eventually reverses. This pattern is consistent with compensatory transfers from migrants in other states.

\paragraph{Related Literature.} The paper is closest to work that studies remittances as a form of household insurance \parencite{yang_choi_2007, bettin_2023}. Most existing evidence studies shocks in receiving countries, where migrants abroad may increase transfers to insure households at origin. We instead study a sender-side shock, where the capacity to remit changes in the destination labor market while origin municipalities are not directly affected. This distinction is important because the same remittance channel that insures households against origin shocks may transmit destination shocks back to those households.

The paper also relates to research on migrant networks and the international transmission of local shocks \parencite{munshi2003, caballero_2023}. The closest paper is \textcite{caballero_2023}, which shows that Mexican households linked to U.S. labor markets hit by the Great Recession became less likely to receive remittances. We differ by studying a localized natural disaster, constructing bilateral state-to-municipality remittance corridors, and focusing on intensive-margin flows and cross-corridor responses.

The structural framework connects to research that models migrants as allocating resources across origin and destination locations \parencite{albert_monras_2022}. The difference is that our model uses this logic to study remittance behavior after a destination-side shock. It also connects to work on interference and spillovers in causal inference \parencite{hudgens2008, aronow2017, butts2023}. In our setting, interference is not only a threat to identification. It is part of the economic mechanism, because migrants linked to the same origin municipality may jointly insure households against shocks that affect only one destination.

The remainder of the paper is organized as follows. Chapter 2 develops the remittance model, describes Hurricane Ian as the natural experiment, and presents the empirical strategy. Chapter 3 describes the construction and validation of the bilateral remittance panel. Chapter 4 reports the estimates for direct effects, exposure effects, and spillovers. Chapter 5 concludes.
